\chapter*{Tóm tắt}
\addcontentsline{toc}{chapter}{Tóm tắt}

Các hệ thống điều phối xe đi chung theo yêu cầu hiện nay đều báo cho hành khách một giờ đón
và một giờ trả dự kiến, rồi sau đó tiếp tục tối ưu lại kế hoạch mỗi khi có khách mới. Hệ quả
là giờ đã báo bị dời đi nhiều lần trong cùng một chuyến, trong khi mọi thước đo chuẩn của
ngành --- tỉ lệ phục vụ, thời gian chờ, độ vòng đường --- đều không ghi nhận điều đó. Một
hành khách được báo đúng một lần và một hành khách bị đổi giờ mười hai lần là giống hệt nhau
trong các hàm mục tiêu của dòng bài điều phối kinh điển, miễn là cả hai vẫn nằm trong cửa sổ
thời gian hợp lệ.

Khóa luận này đặt câu hỏi ngược lại: \emph{một hãng gọi xe chung có thể đưa ra lời hứa về giờ
giấc mà hãng thực sự giữ được hay không, và cái đó đắt bao nhiêu?} Để trả lời, chúng tôi xây
dựng cơ chế \textbf{lời hứa có biên được chứng nhận}. Mỗi lần hệ thống công bố lại một lời
hứa, độ dời được ghi vào một sổ cái riêng của từng hành khách, cộng dồn và chỉ tăng, rồi đối
chiếu với một trần $\beta$; phương án nào làm sổ cái vượt trần sẽ bị loại khỏi tập lựa chọn
ngay tại thời điểm quyết định. Khác biệt cốt lõi so với văn liệu là ở chỗ độ dời được phân rã
theo nguyên nhân: bằng cách chiếu lại chính kế hoạch cũ từ vị trí thật hiện tại của xe, hệ
thống tách được phần do nhà điều hành tự gây ra khỏi phần thay đổi khi giữ nguyên kế hoạch, và
chỉ tính vào sổ cái phần thứ nhất.

Phần lý thuyết đặt bốn họ ràng buộc lời hứa --- theo hành vi của nhà điều hành, hạn chót một
chiều tính từ lời hứa đầu, hạn chót hai chiều, và theo kết quả khách nhìn thấy --- vào cùng một
khung, và chứng minh rằng, xét riêng phần cam kết, chỉ ràng buộc theo hành vi không bao giờ là
nguyên nhân làm rỗng tập phương án khi thế giới trôi dạt; họ nào trong ba họ còn lại hỏng phụ
thuộc chiều trôi. Chứng minh là sơ cấp; giá trị nằm ở khung so sánh và ở đối chiếu với hệ thật.
Phần thực nghiệm gồm 558 lần chạy mô phỏng trên dữ liệu chuyến đi công khai của Manhattan, sử
dụng nền mô phỏng FleetPy và một panel 20 ô đã niêm phong. Trong chế độ không cam kết, phần do
chính bộ điều phối gây ra bằng $89{,}06\%$ lượng dịch chuyển giờ hẹn mà hành khách nhìn thấy; một
trần 30 giây đưa tỉ lệ đó xuống $1{,}37\%$ với cái giá $7{,}13$ điểm phần trăm tỉ lệ phục vụ. Dưới
dữ liệu tắc đường đo thật của Manhattan (thăm dò), cách tính theo kết quả vô nghiệm do cam kết ở
19 trên 20 ô, trong khi cách tính theo hành vi không vô nghiệm do cam kết ở ô nào --- dù 11 ô vẫn
dừng vì một bất biến cửa sổ đón nằm ngoài cơ chế cam kết. Một thí nghiệm thăm dò có tiền đăng ký
kiểm thêm hai họ hạn chót, với dự đoán theo từng ô niêm phong trước khi chạy các nhánh được dự
đoán: dự đoán đúng ở
$98/102$ cặp trên panel chính và ở $31/32$, $40/40$ cặp trên khung sáng và khung tối của một tập giữ
lại; khi tắc đường giảm, hạn chót một chiều vô nghiệm ở 0/8 ô còn hạn chót hai chiều ở 7/8 ô, đúng
chiều định lý dự đoán.

\vspace{6pt}
\noindent\textbf{Từ khóa:} điều phối xe đi chung, lời hứa có biên, độ ổn định lịch trình,
quy trách nguyên nhân, nghĩa vụ phương tiện, đánh đổi phục vụ--ổn định.

\clearpage

\chapter*{Abstract}
\addcontentsline{toc}{chapter}{Abstract}

Modern on-demand ride-pooling dispatchers announce an estimated pickup and drop-off time to
every rider, then keep re-optimising the plan as new requests arrive. As a result, the
announced times are revised repeatedly within a single trip, while none of the field's
standard metrics --- service rate, waiting time, detour --- record this at all. A rider told
once and a rider re-quoted twelve times are indistinguishable under the objective functions of
the classical dispatching literature, as long as both remain inside the admissible time window.

This thesis asks the opposite question: \emph{can a ride-pooling operator make a timing
promise it can actually keep, and what does that cost?} We build a mechanism of
\textbf{certified bounded promises}. Every time the system republishes a promise, the
displacement is recorded in a per-rider ledger that accumulates monotonically and is checked
against a cap $\beta$; any candidate plan that would push the ledger past the cap is removed
from the choice set at decision time. The core departure from the literature is that the
displacement is decomposed by cause: by re-projecting the unchanged previous plan from the
vehicle's true current position, the system separates what the operator caused from what
changes when the plan is kept, and charges only the former.

The theoretical part places four families of promise constraints --- a cap on the operator's
conduct, a one-sided deadline measured from the first promise, a two-sided deadline, and a cap
on the outcome the rider observes --- in one framework, and proves that, as far as the
commitment gate is concerned, only the conduct cap is never the cause of an empty choice set
when the world drifts; which of the other three fails depends on the direction of the drift.
The proof is elementary; its value lies in the common frame and in checking it against a real
system. The experimental part comprises 558 simulation runs on public Manhattan trip data,
using the FleetPy simulator and a sealed 20-cell panel. Under the unconstrained policy the
operator-induced part equals $89.06\%$ of the rider-visible promise movement; a 30-second cap
brings this share down to $1.37\%$ at a cost of $7.13$ percentage points of service rate. Under
measured Manhattan congestion (exploratory), the outcome basis is commitment-infeasible on 19
of 20 cells while the conduct basis is commitment-infeasible on none --- although 11 cells still
stop on a pickup-window invariant that lies outside the commitment mechanism. An exploratory
preregistered experiment adds the two deadline families, with per-cell predictions sealed before
any run of a predicted arm: they hold on $98/102$ pairs on the main panel and on $31/32$ and $40/40$ pairs in the
morning and evening windows of a held-out set; when congestion eases, the one-sided deadline is
commitment-infeasible on 0 of 8 cells and the two-sided one on 7 of 8, the direction the theorem
predicts.

\vspace{6pt}
\noindent\textbf{Keywords:} ride-pooling dispatch, bounded promises, schedule stability, causal
attribution, obligation of means, service--stability trade-off.

\clearpage
